
\setcounter{footnote}{-1}
\chapter*{Expos\'e \texorpdfstring{VI\textsubscript{A}}{VI_A}\\
Generalities on Algebraic Groups}
\phantomsection
\label{sga3:VIA:title}
\addcontentsline{toc}{chapter}{Expos\'e VI-A. Generalities on Algebraic Groups}
\markboth{EXPOS\'E VI\textsubscript{A}. GENERALITIES ON ALGEBRAIC GROUPS}%
  {EXPOS\'E VI\textsubscript{A}. GENERALITIES ON ALGEBRAIC GROUPS}

\begin{center}
  \textit{by P. Gabriel}
\end{center}

\footnote{Editor's note: version of October 13, 2024.}Throughout this
chapter, \(A\) will denote an Artinian local ring with residue field \(k\).
A group scheme \(G\) over \(\operatorname{Spec} A\) will be called simply
an \(A\)-\emph{group}. This \(A\)-group is said to be \emph{locally of
finite type} if its underlying scheme is locally of finite type over \(A\);
it is said to be \emph{algebraic} if its underlying scheme is of finite
type over \(A\).

\section*{0. Preliminary Remarks}
\phantomsection
\label{sga3:VIA:0}
\addcontentsline{toc}{section}{0. Preliminary Remarks}

\medskip
\noindent\textbf{0.1.}
\phantomsection
\label{sga3:VIA:0.1}
Let us first consider a group scheme \(G\) over an arbitrary scheme \(S\).
We call the structural morphism
\[
  \mu:G\times_S G\longrightarrow G
\]
\emph{multiplication}, and we call the morphism \(c:G\to G\), defined by
\(c(T)(x)=x^{-1}\) for every scheme \(T\) over \(S\) and every
\(x\in G(T)\), \emph{inversion}. If \(U\) and \(V\) are subsets of the
underlying set of \(G\), we denote by \(U\cdot V\) the image under
multiplication of the subset of \(G\times_S G\) consisting of those points
whose first projection belongs to \(U\) and whose second projection belongs
to \(V\). Likewise, the notations \(U^{-1}\) and \(c(U)\) are equivalent.

Let \(\operatorname{pr}_1\) be the projection of \(G\times_S G\) onto the
first factor, and let
\[
  \sigma:G\times_S G\longrightarrow G\times_S G
\]
be the morphism with components \(\operatorname{pr}_1\) and \(\mu\). For
every \(S\)-scheme \(T\), the map \(\sigma(T)\) is
\[
  (x,y)\longmapsto (x,xy);
\]
it follows that \(\sigma\) is an automorphism. The composite of this
automorphism with the projection \(\operatorname{pr}_2\) of
\(G\times_S G\) onto the second factor is the multiplication morphism.
When \(G\) is flat over \(S\), \(\operatorname{pr}_2\), and hence \(\mu\),
are flat morphisms; when \(G\) is smooth over \(S\),
\(\operatorname{pr}_2\), and hence \(\mu\), are smooth morphisms, and so
on.

\medskip
\noindent\textbf{0.2.}
\phantomsection
\label{sga3:VIA:0.2}
From now on we assume that \(S\) is the spectrum of an Artinian local ring
\(A\) with residue field \(k\). We denote by
\((\mathbf{Sch}_{/k})_{\mathrm{red}}\) the category of reduced schemes over
\(k\). For every scheme \(X\) over \(A\), the reduced scheme
\(X_{\mathrm{red}}\) is an object of
\((\mathbf{Sch}_{/k})_{\mathrm{red}}\), and the functor
\[
  X\longmapsto X_{\mathrm{red}}
\]
is right adjoint to the inclusion of
\((\mathbf{Sch}_{/k})_{\mathrm{red}}\) into \(\mathbf{Sch}_{/A}\).
It follows that, for every \(A\)-group \(G\), \(G_{\mathrm{red}}\) is a
group in the category
\((\mathbf{Sch}_{/k})_{\mathrm{red}}\).\footnote{Editor's note: the
sentences that follow have been added.} In other words, for every reduced
\(k\)-scheme \(T\), \(G_{\mathrm{red}}(T)\) is endowed, functorially in
\(T\), with a group structure. One must bear in mind that
  \(G_{\mathrm{red}}\) is not necessarily a \(k\)-group, because its
  ``multiplication'' is only a morphism
\stepcounter{footnote}%
\[
  (G_{\mathrm{red}}\times_k G_{\mathrm{red}})_{\mathrm{red}}
  \longrightarrow
  G_{\mathrm{red}}\text{\hyperlink{sga3:VIA:fn:red-examples}{%
    \textsuperscript{\normalfont\thefootnote}}}.
\]
\footnotetext[\value{footnote}]{\hypertarget{sga3:VIA:fn:red-examples}{}Editor's note: for examples of group schemes \(G\) over a
nonperfect field \(k\) for which \(G_{\mathrm{red}}\) is not a group scheme
over \(k\), see \SGARefLink{ex:sga3-VIA-1.3.2}{1.3.2} below.}

\footnote{Editor's note: the remainder of \SGARefLink{sga3:VIA:0.2}{0.2}, as well as \SGARefLink{sga3:VIA:0.3}{0.3}, has been
slightly modified.}
However, if \(k\) is a perfect field, the inclusion of
\((\mathbf{Sch}_{/k})_{\mathrm{red}}\) into \(\mathbf{Sch}_{/k}\)
commutes with products. Groups in
\((\mathbf{Sch}_{/k})_{\mathrm{red}}\) may therefore be identified with
group schemes over \(k\) whose underlying schemes are reduced. In this
case, if \(G\) is a \(k\)-group, \(G_{\mathrm{red}}\) is a group subscheme
of \(G\); in general, this subgroup is not normal in \(G\).

For example, if \(k\) is a field of characteristic \(3\), the constant
group \((\mathbb Z/2\mathbb Z)_k\) acts nontrivially on the diagonalizable
group \(D_k(\mathbb Z/3\mathbb Z)\) (cf. \SGARefLink{sga3:i:subsection:4-1}{Expos\'e I, 4.1} and \SGARefLink{cor:sga3-VIB-4.4}{4.4}). If
\(G\) denotes the semidirect product of \(D_k(\mathbb Z/3\mathbb Z)\) by
\((\mathbb Z/2\mathbb Z)_k\) defined by this action, then
\(G_{\mathrm{red}}\) is identified with \((\mathbb Z/2\mathbb Z)_k\) and
is not normal in \(G\).\footnote{Editor's note: for an analogous example
with \(G\) connected, one may take, when
\(\operatorname{char}(k)=p>0\), the semidirect product of
\(\alpha_{p,k}\) and \(\mathbb G_{m,k}\).}

Let \(k\) be an arbitrary field, let \(k^{p^{-\infty}}\) be its perfect
closure, and let \(H\) be a group in
\((\mathbf{Sch}_{/k})_{\mathrm{red}}\). Then
\[
  (H\otimes_k k^{p^{-\infty}})_{\mathrm{red}}
\]
is a group scheme over the perfect field \(k^{p^{-\infty}}\). Since
\(H\otimes_k k^{p^{-\infty}}\) and
\((H\otimes_k k^{p^{-\infty}})_{\mathrm{red}}\) have the same underlying
topological space, groups in \((\mathbf{Sch}_{/k})_{\mathrm{red}}\) share
with \(k\)-groups certain topological properties that are invariant under
extension of the base field. For example, it will follow from \hyperref[sga3:VIA:0.3]{0.3} and the
preceding remarks that every group in
\((\mathbf{Sch}_{/k})_{\mathrm{red}}\) is separated.

In what follows, we shall encounter groups in
\((\mathbf{Sch}_{/k})_{\mathrm{red}}\) whenever we deal with a nonempty
locally closed subset \(U\) of an \(A\)-group \(G\) such that
\[
  U\cdot U=U
  \qquad\text{and}\qquad
  U^{-1}=U.
\]
Indeed, the reduced subscheme of \(G\) defined by \(U\) is then a group in
\((\mathbf{Sch}_{/k})_{\mathrm{red}}\).

\medskip
\noindent\textbf{0.3.}
\phantomsection
\label{sga3:VIA:0.3}
\textit{An \(A\)-group \(G\) is always separated, because the unit section
\(e:\operatorname{Spec}A\to G\) is a closed immersion.}
Indeed, let \(x\) be the unique point of \(\operatorname{Spec}A\), and let
\(\eta:G\to\operatorname{Spec}A\) be the structural morphism. Since
\(\eta\circ e=\operatorname{id}_{\operatorname{Spec}A}\), for every
affine open \(U=\operatorname{Spec}B\) of \(G\) containing \(e(x)\), the
homomorphism \(B\to A\) has a section and is therefore surjective. It
follows that \(e\) is a closed immersion.\footnote{Editor's note: this
argument is valid for every local ring \(A\) of dimension zero, and shows
that, if \(S\) is a discrete scheme, every \(S\)-group is separated; cf.
VI\textsubscript{B}, \SGARefLink{prop:sga3-VIB-5.2}{5.2}. On the other hand (cf.
VI\textsubscript{B}, \SGARefLink{ctr:sga3-VIB-5.6.4}{5.6.4}), if \(S\) contains a nonisolated closed point
\(s\), the \(S\)-scheme \(G\) obtained by gluing two copies of \(S\) along
the open subset \(S-\{s\}\) is not separated over \(S\), although it
carries an \(S\)-group structure.}

The diagonal in \(G\times_A G\) is identified with the functor from
\((\mathbf{Sch}_{/A})^\circ\) to \(\mathbf{Sets}\) which associates to
every \(A\)-scheme \(S\) the inverse image of the identity element of
\(G(S)\) under the map
\[
  \varphi(S):G(S)\times G(S)\longrightarrow G(S),
  \qquad (x,y)\longmapsto x\cdot y^{-1}.
\]
We therefore have the following cartesian square. Since the diagonal
morphism is obtained from a closed immersion by base change, it is itself
a closed immersion.

\begin{center}
  \phantomsection
  \label{sga3:VIA:0.3-cartesian-square}
  \begin{tikzcd}[row sep=large,column sep=large]
    G\times_S G
      \arrow[r,"\varphi"]
    &
    G
    \\
    G
      \arrow[u,"\text{diagonal morphism}"]
      \arrow[r]
    &
    \operatorname{Spec}A
      \arrow[u,"\text{unit section}"']
  \end{tikzcd}
\end{center}


\medskip
\noindent\textbf{0.4.}\footnote{Editor's note: the original has been
expanded in what follows; in particular, \hyperref[sga3:VIA:0.4.1]{Remark 0.4.1} has been added.}
\phantomsection
\label{sga3:VIA:0.4}
Let \(G\) be an \(A\)-scheme. We shall say that a point \(g\) of \(G\) is
\emph{strictly rational over \(A\)} if there exists an \(A\)-morphism
\[
  s:\operatorname{Spec}A\longrightarrow G
\]
which sends the unique point of \(\operatorname{Spec}A\) to \(g\); that
is, if the homomorphism \(A\to\mathcal O_{G,g}\) admits a retraction. One
then has \(\kappa(g)=k\), so \(g\) is a closed point of \(G\). Indeed, if
\(B\) is the ring of an affine open neighborhood of \(g\), and if
\(\mathfrak p\) is the prime ideal of \(B\) corresponding to \(g\), then
\(B/\mathfrak p\subset k\) is a finite \(A\)-algebra, hence an integral
Artinian ring, and therefore a field.

Suppose henceforth that \(G\) is an \(A\)-group. Such a morphism
\(s:\operatorname{Spec}A\to G\) defines an automorphism \(r_s\) of the
scheme \(G\) over \(A\), called \emph{right translation by \(s\)}. For
every morphism \(\pi:S\to\operatorname{Spec}A\), the automorphism
\(r_s(\pi)\) of \(G(S)\) is defined by
\[
  x\longmapsto x\cdot G(\pi)(s),
  \qquad x\in G(S).
\]
Similarly, \(\ell_s\) denotes \emph{left translation by \(s\)}, namely
the automorphism of \(G\) defined by
\[
  \ell_s(\pi)(x)=G(\pi)(s)\cdot x,
  \qquad x\in G(S).
\]

The schemes \(G\otimes_A k\) and \(G\) have the same underlying
topological space \(\underline G\), \(G\otimes_A k\) is a \(k\)-group,
and \(s\otimes_A k\) depends only on \(g\), not on \(s\). Consequently,
the automorphisms of \(\underline G\) induced by \(r_s\) and \(\ell_s\)
(or by \(r_{s\otimes k}\) and \(\ell_{s\otimes k}\)) depend only on \(g\),
not on \(s\). Thus, if \(P\) is a subset of \(\underline G\), we may write
\[
  r_g(P)\ \text{or}\ P\cdot g
  \qquad
  \bigl(\text{respectively, }\ell_g(P)\ \text{or}\ g\cdot P\bigr)
\]
instead of \(r_s(P)\) (respectively, \(\ell_s(P)\)), in agreement with
\hyperref[sga3:VIA:0.1]{0.1}.

\medskip
\noindent\textbf{Remark 0.4.1.}
\phantomsection
\label{sga3:VIA:0.4.1}
\textit{If \(g\) is a strictly rational point of \(G\), and if
\(A\to A'\) is a morphism of Artinian local rings, then
\(G'=G\otimes_A A'\) has a unique point \(g'\) above \(g\), and \(g'\) is
strictly rational over \(A'\). Moreover, if \(P'\) denotes the inverse
image of \(P\) in \(G'\), then \(P'\cdot g'\) is the inverse image of
\(P\cdot g\) (cf. EGA I, 3.4.8).}

\medskip
\noindent\SGAInlineRefTarget{sga3:vi:proposition:0-5}\textbf{Proposition 0.5.}\footnote{Editor's note: the original
stated this result under the hypothesis that \(G\) be locally of finite
type over \(A\). Since the result is useful in the general case, and
since the proof is essentially the same, it has been stated and proved
here in full generality. It will be used several times below.}
\phantomsection
\label{sga3:VIA:0.5}
\textit{Let \(G\) be an \(A\)-group, and let \(U,V\) be two dense open
subsets of \(G\). Then \(U\cdot V\), that is, the image of
\(U\times_A V\) under multiplication, is the entire underlying space of
\(G\).}

\begin{proof}
Since \(G\) and \(G\otimes_A k\) have the same underlying space, after
replacing \(A\) by \(k\) and \(G\) by \(G\otimes_A k\), if necessary, we
may assume that \(A=k\). Let \(g\in G\), and set \(K=\kappa(g)\). Then
left translation \(\ell_g\) is an automorphism of \(G_K\). Since the
projection \(G_K\to G\) is open, \(U_K\) and \(V_K\) are dense open
subsets of \(G_K\), as is the image of \(V_K\) under \(\ell_g\).
Consequently there exists \(v\in V_K\) such that
\(u=\ell_g(v)\) belongs to \(U_K\).

Let \(L\) be an extension of \(K\) containing \(\kappa(v)\), and hence
\(\kappa(u)\), and let \(g_L\) and \(v_L\) be the \(L\)-points of \(G_L\)
induced by \(g\) and \(v\). Then
\[
  g_L\cdot v_L=u'
\]
is a point of \(G_L\) above \(u\), and therefore
\[
  g_L=u'\cdot v_L^{-1}
\]
lies above \(U\cdot V\). Thus \(g\in U\cdot V\), proving the proposition.
\end{proof}

\medskip
\noindent\textbf{Corollary 0.5.1.}\footnote{Editor's note: this corollary
has been inserted here; cf. \SGARefLink{prop:sga3-VIA-2.4}{2.4} and \SGARefLink{cor:sga3-VIA-2.6.2}{2.6.2} below.}
\phantomsection
\label{sga3:VIA:0.5.1}
\textit{If \(G\) is an irreducible \(A\)-group, then \(G\) is
quasi-compact.}

\begin{proof}
Let \(U\) be a nonempty affine open subset of \(G\). Then \(U\) is dense
in \(G\), so by \hyperref[sga3:VIA:0.5]{0.5} the morphism
\[
  \mu:U\times_A U\longrightarrow G
\]
is surjective. Hence \(G\) is quasi-compact, since \(U\times_A U\) is.
\end{proof}

\medskip
\noindent\textbf{Corollary 0.5.2.}\footnote{Editor's note: this
corollary, used implicitly in VIII, \SGARefLink{sga3:viii:corollary:6:7}{6.7}, has been added; see also
VI\textsubscript{B}, \SGARefLink{cor:sga3-VIB-6.2.5}{6.2.5}.}
\phantomsection
\label{sga3:VIA:0.5.2}
\textit{Let \(G\) be an \(A\)-group scheme, and let \(H\) be a group
subscheme of \(G\) over \(A\). Then \(H\) is closed.}

\begin{proof}
Let \(k'\) be the perfect closure of the residue field \(k\) of \(A\).
Since the underlying topological spaces of \(G\) and \(H\) are unchanged
under the base change \(A\to k\to k'\), we may assume that \(A=k\) is a
perfect field. We may then assume that \(G\) and \(H\) are reduced, hence
geometrically reduced.

Let \(\overline H\) be the closure of \(H\). Then
\(\mu^{-1}(\overline H)\) is a closed subset of \(G\times G\) containing
\(H\times H\). The morphisms
\[
  H\longrightarrow\operatorname{Spec}k
  \qquad\text{and}\qquad
  \overline H\longrightarrow\operatorname{Spec}k
\]
are universally open. Since \(H\) is dense in \(\overline H\),
\(H\times H\) is dense in \(\overline H\times H\), and
\(\overline H\times H\) is dense in
\(\overline H\times\overline H\). Thus \(H\times H\) is dense in
\(\overline H\times\overline H\). Consequently
\[
  \mu(\overline H\times\overline H)\subset\overline H.
\]
Since \(\overline H\times\overline H\) is reduced, \(\mu\) induces a
morphism
\[
  \mu':\overline H\times\overline H\longrightarrow\overline H.
\]

Now let \(g\in\overline H\), and put \(K=\kappa(g)\). The projection
\(\overline H_K\to\overline H\) is open, so \(H_K\) and
\(\ell_g(H_K)\) are two dense open subsets of \(\overline H_K\). Hence
there exist \(u,v\in H_K\) such that \(\ell_g(v)=u\). As in the proof of
\SGARefLink{sga3:vi:proposition:0-5}{0.5}, it follows that \(g\) belongs to \(H\cdot H=H\). Therefore
\(\overline H=H\).
\end{proof}

